% SOURCE_AUTHORITY: sga5_fr_workpass.tex, lines 8006--8032 (LF-normalized unit).
% SOURCE_SHA256: 791F4EFFC5E02832D5D77ED03518C8156D6F07E4C8238B03545DB93D883FBB28
c) Quando \(X\) verifica a condição de Mittag-Leffler, os diferentes sistemas projetivos \(\gr^p(X)\) são essencialmente constantes (EGA \(0_{\mathrm{III}}\) 13.4.3). Daí segue imediatamente que, para \(r<+\infty\), os sistemas projetivos \(E_r^{pq}(X)\) são essencialmente constantes. Denotaremos por
\[
\mathcal E_r^{pq}(X)
\]
seus limites projetivos. Ao passar ao limite os operadores diferenciais, obtém-se assim um filtro espectral \(\mathcal E_r^{pq}(X)\) \((p,q\in\mathbb Z;\ r\ge1)\). Os objetos
\[
\mathcal Z_r^{pq}(X),\quad \mathcal B_r^{pq}(X)\quad(r<+\infty)
\]
associados a este filtro espectral são precisamente os limites dos sistemas projetivos essencialmente constantes \(Z_r^{pq}(X)\) e \(B_r^{pq}(X)\).

Da parte b) segue também que
\[
B_\infty^{pq}(X)=B_{p+1}^{pq}(X),
\]
que permite definir
\[
\mathcal B_\infty^{pq}(X)=\mathcal B_{p+1}^{pq}(X).
\]
Supondo ainda que \(X\) verifique (ML), como, para todo \(\alpha\),
\[
B_r^{pq}(X_\alpha)=B_{r+1}^{pq}(X_\alpha)\qquad(r>p),
\]
tem-se uma sucessão de monomorfismos canônicos:
\begin{equation}
\cdots\subset\mathcal E_r^{pq}(X)\subset\cdots\mathcal E_{p+2}^{pq}(X)\subset\mathcal E_{p+1}^{pq}(X).
\tag{S}
\end{equation}
